% Part: many-valued-logic
% Chapter: syntax-and-semantics
% Section: formulas

\documentclass[../../../include/open-logic-section]{subfiles}

\begin{document}

\olfileid{mvl}{syn}{fml}

\olsection{\usetoken{P}{formula}}

\begin{defn}[સૂત્ર]
\ollabel{defn:formulas}
વિધાનાત્મક ભાષા~$\Lang L$નાં \emph{!!{formula}s}નો
ગણ~$\Frm[L]$ નીચે મુજબ અનુમાનાત્મક રીતે વ્યાખ્યાયિત થાય છે:
\begin{enumerate}
\item દરેક !!{propositional variable}~$\Obj p_i$ આણ્વિક
  !!{formula} છે.
\item $\Lang L$નો દરેક $0$-સ્થાનીય સંયોજક
  (વિધાનાત્મક અચલ) આણ્વિક !!{formula} છે.
\item $\star$ એ $\Lang L$નો $n$-સ્થાનીય સંયોજક
  હોય અને $!A_1$, \dots, $!A_n$ !!{formula}s હોય,
  તો $\star(!A_1, \dots, !A_n)$ પણ !!a{formula} છે.
\tagitem{limitClause}{આ સિવાય બીજું કંઈ !!a{formula} નથી.}{}
\end{enumerate}
$\star$ $1$-સ્થાનીય હોય તો $\star(!A_1)$ને ઘણી વાર
માત્ર $\star !A_1$ લખીશું. $\star$ $2$-સ્થાનીય હોય તો
$\star(!A_1,!A_2)$ને ઘણી વાર $(!A_1 \star !A_2)$
લખીશું.
\end{defn}

હંમેશની જેમ, ઘણી વાર સૌથી બહારનાં કૌંસ ચૂપચાપ
છોડી દઈશું.

\begin{ex}
  પ્રમાણભૂત ભાષા~$\Lang{L_0}$માં $\Obj p_1 \lif (\Obj p_1
  \land \lnot \Obj p_2)$ એક સૂત્ર છે. ગુણન
  તર્કશાસ્ત્રની ભાષામાં તેને $\Obj p_1 \lif (\Obj p_1 \odot
  \lnot \Obj p_2)$ તરીકે લખાશે. ભાષામાં $1$-સ્થાનીય
  $\triangle$ ઉમેરીએ તો $\triangle (\Obj p_1
  \land \Obj p_2) \lif (\triangle \Obj p_1 \land \triangle \Obj p_2)$
  જેવાં સૂત્રો પણ મળશે.
\end{ex}

\end{document}
